\begingroup
\setlength{\emergencystretch}{2em}
\hypertarget{source-membership-results-and-limitations}{%
\subsection{Source-membership results and
limitations}\label{source-membership-results-and-limitations}}

\hypertarget{fixed-development-performance}{%
\subsubsection{Fixed development
performance}\label{fixed-development-performance}}

In the frozen 234-address test partition, the tree recovers 97 of 98
core memberships and assigns one of 136 outside-core comparisons to the
core class. Membership precision, recall and F1 are each 0.989796. These
metrics evaluate the operational membership labels, not independently
verified theft or hardware identity.

\begin{longtable}[]{@{}llllll@{}}
\toprule\noalign{}
Partition & AP & ROC-AUC & Precision & Recall & F1 \\
\midrule\noalign{}
\endhead
\bottomrule\noalign{}
\endlastfoot
Validation & 0.981534 & 0.990422 & 0.982456 & 1.000000 & 0.991150 \\
Test & 0.993757 & 0.994748 & 0.989796 & 0.989796 & 0.989796 \\
\end{longtable}

\begin{longtable}[]{@{}lllll@{}}
\toprule\noalign{}
Partition & TP & FP & FN & TN \\
\midrule\noalign{}
\endhead
\bottomrule\noalign{}
\endlastfoot
Train & 346 & 6 & 0 & 681 \\
Validation & 56 & 1 & 0 & 54 \\
Test & 97 & 1 & 1 & 135 \\
\end{longtable}

Train counts are resubstitution diagnostics. Validation and test are
held out from fitting, but belong to the same previously inspected
incident/window; they are not an untouched or independent-incident
evaluation. All AP and ROC-AUC values were independently checked from
tied score groups and pairwise ranking counts. Confusion counts and F1
also reproduce directly from the saved predictions.

\begin{longtable}[]{@{}lll@{}}
\toprule\noalign{}
Fixed test method & Membership AP & Membership ROC-AUC \\
\midrule\noalign{}
\endhead
\bottomrule\noalign{}
\endlastfoot
Training-prevalence constant & 0.418803 & 0.500000 \\
Longer inactivity & 0.982880 & 0.986645 \\
Larger balance & 0.765014 & 0.911877 \\
Prespecified tree & 0.993757 & 0.994748 \\
\end{longtable}

Inactivity alone already achieves test AP 0.982880. The tree's
descriptive gain is 1.09 percentage points on this fixed partition. No
independent significance or causal claim follows, and no additional
parameter search or model was fitted after inspecting these results.

\hypertarget{rules-and-errors}{%
\subsubsection{Rules and errors}\label{rules-and-errors}}

The fitted tree has five leaves and uses only prior balance and
inactivity. At the fixed decision threshold, the positive-membership
region requires float32-converted balance greater than 14,820,511.5
satoshis and inactivity greater than 1,006.5 blocks. A further split at
8,750 blocks distinguishes two positive leaf scores. This is the fitted
cohort rule, not a proposed operational theft/device threshold. A
redundant split between two zero-score comparison leaves is retained
without post-hoc pruning. Exact replay conditions and per-partition
node/leaf supports remain in the saved model and rules.

No explicit missingness indicator was selected, and no observed source
reached a numeric split using an imputed value for that split variable.
The full training-only preparation remains exported. This observation
does not characterize the unknown vectors of excluded histories.

\begin{longtable}[]{@{}
  >{\raggedright\arraybackslash}p{(\columnwidth - 6\tabcolsep) * \real{0.2500}}
  >{\raggedright\arraybackslash}p{(\columnwidth - 6\tabcolsep) * \real{0.2500}}
  >{\raggedright\arraybackslash}p{(\columnwidth - 6\tabcolsep) * \real{0.2500}}
  >{\raggedright\arraybackslash}p{(\columnwidth - 6\tabcolsep) * \real{0.2500}}@{}}
\toprule\noalign{}
\begin{minipage}[b]{\linewidth}\raggedright
Partition / membership error
\end{minipage} & \begin{minipage}[b]{\linewidth}\raggedright
Source ID
\end{minipage} & \begin{minipage}[b]{\linewidth}\raggedright
Prior balance (sat)
\end{minipage} & \begin{minipage}[b]{\linewidth}\raggedright
Inactivity (blocks)
\end{minipage} \\
\midrule\noalign{}
\endhead
\bottomrule\noalign{}
\endlastfoot
validation / FP & \texttt{comparison\_\allowbreak{}0328} & 16604900 & 145611 \\
test / FN & \texttt{cohort\_\allowbreak{}0191} & 69145962 & 409 \\
test / FP & \texttt{comparison\_\allowbreak{}0401} & 40096077 & 7000 \\
\end{longtable}

The missed case source \texttt{cohort\_\allowbreak{}0191} has 105 earlier receipt
transactions and last activity only 409 blocks before reference. It
falls outside the learned inactivity region despite its case membership.
The test comparison \texttt{comparison\_\allowbreak{}0401} has one earlier receipt, a
balance of 40,096,077 satoshis and 7,000 blocks of inactivity,
resembling the covered case profile. The validation error is also an
old-funded comparison. These are membership errors; neither predictions
nor corrections settle compromise or device status.

\hypertarget{coverage-and-inference-limits}{%
\subsubsection{Coverage and inference
limits}\label{coverage-and-inference-limits}}

Coverage is 100\% for cases and 87.8\% for selected comparisons, ranging
from 81.36\% to 91.51\% across comparison strata. The sampled core
prevalence increases from 33.33\% of selections to 36.28\% of available
vectors. The 122 unavailable comparisons have median saved
full-current-history transaction count 3,347, versus 3 for available
comparisons. These counts were observed at different collection
snapshots and are diagnostic metadata, never prior-block predictors or
estimates of missing features. Missingness cannot be assumed random.

All 500 positives come from one retrospectively reconstructed incident.
Core membership and case evidence were learned after the event; their
publication/availability dates are not established for prospective use.
Addresses are not independent victims or devices. Same-chain
reconstruction agreement, an owner report, provider coverage and genuine
device/compromise verification are distinct evidence tiers.
Independently verified compromise/device labels remain unknown for every
selection, and comparisons are not verified benign examples.

The result supports a reproducible, interpretable description of covered
source profiles and a narrow within-window membership baseline. It does
not establish a cause of loss, a signing-device fingerprint, calibrated
theft probabilities, independent-incident transfer or general blockchain
detection performance. The selection prevalence is study-defined. The
122 bounded-incomplete prefixes remain unresolved; a validated prefix
also does not certify a complete current history. The frozen external
nineteen-feature forest and all earlier experiment artifacts remain
unchanged.

\hypertarget{reproducibility-and-supported-conclusion}{%
\subsubsection{Reproducibility and supported
conclusion}\label{reproducibility-and-supported-conclusion}}

The fitted experiment is pinned to
\href{https://github.com/HnsHugo/STAGE-UFU/commit/00659df35d3140143fa12b37ba35ea56722a4854}{commit
00659df} on branch \texttt{task-\allowbreak{}2-\allowbreak{}wallet-\allowbreak{}analysis}. Its protocol freezes
the data, component assignments, preparation, tree parameters, threshold
and software versions. All 1,378 JSON-exported scores exactly match the
fitted estimator. Thirty focused tests and verification without a model
refit passed in the saved execution. The report builder checks the saved
predictions, confusion counts, F1, fingerprints, exclusions and
component partitioning; it performs no new training or collection.

A separate independent evaluation would require a justified additional
incident/source universe, evidence-backed label definitions and a new
prespecified protocol. The current validation/test outcomes must not be
reused for tuning or described as newly independent evidence.

\endgroup
